\cotefile{81}
\scanpage{19}
Dém. par récurrence, c'est trivial si $\operatorname{card} I=0,1$ Le cas $\operatorname{card} I=2$ a été vu. Cas général \unsure{avec} $\operatorname{card} I\geq 3$, on choisit $J$ comme dessus. Il suffit de voir que les \unsure{ens.}
\[
\unsure{U|\deleted{=}U}_{J,i}\cap\Sigma_i^+,\ U_{J,i}\cap\Sigma_i^-
\]
ont la forme précédente, il suffit de le voir pour \unsure{l'une|l'un} des deux.

Par hyp. de récurrence, \unsure{on a}
\[
U_{J,i}=X\setminus\mathop{\unsure{\coprod|\bigsqcup}}\limits_{\alpha\in \deleted{\unsure{J}}K_0}\widehat{\Sigma}_\alpha^+
\]
où
\[
\left\{\begin{aligned}
&\deleted{\Sigma_{\unsure{\alpha|0}}}K_0\subset J\\
&\Sigma_\alpha^+\in\pi_0(X\setminus E_\alpha)\qquad (\forall\,\alpha\in K_0)\\
&\widehat{\Sigma}_\alpha^+=\Sigma_\alpha^+\cup E_\alpha.
\end{aligned}\right.
\]
\unsure{On a} donc
\[
U=U_{J,i}\cap\Sigma_i^-=X\setminus\left(\coprod_{\alpha\in K_0}\widehat{\Sigma}_\alpha^+\ \cup\ \widehat{\Sigma}_i^+\right)
\]
\deleted{Soit} Au deuxième $=$, dans la \deleted{\unsure{somme|réunion}} réunion $\bigcup_{\alpha\in K_0}\widehat{\Sigma}_\alpha^+$, il suffit de se borner à \unsure{sommer} sur les \deleted{$\Sigma$} $\alpha\in K_0$ tels que $\widehat{\Sigma}_\alpha^+\not\subset\widehat{\Sigma}_i^+$ i.e. $\Sigma_\alpha^+\not\subset\Sigma_i^+$ soit $K_1$ \unsure{cet} ens., donc
\[
U=X\setminus\left(\bigcup_{\alpha\in K_1}\widehat{\Sigma}_\alpha^+\ \cup\ \widehat{\Sigma}_i^+\right)
\]
% archive sheet number: 18
\scanpage{20}
Ainsi les classes de $\Pi_I^{\unsure{\circ|0}}$ sont les classes de $\Pi_J^\circ$ distinctes de $\deleted{\Pi}U_{J,i}$, et de plus les \unsure{clas.|deux} parmi
\[
U_{J,i}\cap\Sigma_{\unsure{\deleted{J}i|j|i}}^+,\ U_{J,i}\cap\Sigma_i^-
\]
qui sont non vides. Or les deux classes en question sont non vides, car \deleted{$\unsure{\Sigma_i|\Sigma_j}$} $U_{J,i}$ est un voisinage de $\deleted{\Sigma}E_{\unsure{i|j}}$, et
\[
\overline{\Sigma_i^+}\cap\deleted{\Sigma}E_i\ne\varnothing,\quad
\overline{\deleted{E}\Sigma_i^-}\cap\deleted{\Sigma}E_i\ne\varnothing.
\]
On \unsure{voit|sait} donc que le nb des classes est en \deleted{\unsure{tout|est}} $(n-1)+2=n+1=\operatorname{card} I+1$ (ce qui achève \unsure{aussi|ainsi} la \unsure{dém.|dem.} par récurrence sur le nb des classes).

\uline{Lemme} Toute classe de $\Pi_I^\circ$ est de la forme $X\setminus\bigcup_{\alpha\in\deleted{\unsure{I}}K}(\overbrace{\Sigma_\alpha^+\cup E_\alpha}^{\inserted{\widehat{\Sigma}_\alpha^+}})$ où $K$ est une partie \inserted{\deleted{non vide}} de $I$, et
% slanted
\reversemarginpar\margin{\makebox[\linewidth][r]{\rotatebox{19}{\shortstack[l]{$K$ est $\ne\varnothing$\\si $I\ne\varnothing$}}}}\normalmarginpar
$\forall\alpha\in K$, \deleted{on a} $\Sigma_\alpha^+\deleted{\unsure{\Rightarrow}}\in\pi_0(X\setminus E_\alpha)$, tels que $\forall\alpha,\beta\in K$, $\alpha\ne\beta$, on ait $\widehat{\Sigma}_\alpha^+\cap\widehat{\Sigma}_\beta^+=\varnothing$ $(\Longleftrightarrow\Sigma_\alpha^+\cap\Sigma_\beta^+=\varnothing)$.
% archive sheet number: 19
\scanpage{21}
Soit $\alpha_0\in K_1$, je dis que \unsure{l'on a|on a}
\[
\widehat{\Sigma}_{\alpha_0}\cap\widehat{\Sigma}_i^+=\varnothing
\qquad\text{\unsure{i.e|i.e.}}\qquad
\Sigma_{\alpha_0}\cap\Sigma_i=\varnothing
\]
\unsure{(ce|Ce} qui achèvera la démonstration, car on aura
\[
U=X\setminus\bigcup_{\alpha\in K}\widehat{\Sigma}_\alpha^+
\qquad\left\{\begin{aligned}
&K=K_1\cup\{i\}\\
&\deleted{\unsure{\Sigma_\alpha}}
\end{aligned}\right.
\]
comme dans le lemme.

Or \unsure{comme} on a
\[
E_i\subset U_{J,i}=X\setminus\mathop{\unsure{\rlap{\raisebox{.65ex}{$\scriptstyle-$}}\bigcup|\bigcup}}\limits_{\alpha\in K_0}\widehat{\Sigma}_\alpha^+
\]
cela implique que
\[
E_i\subset X\setminus\widehat{\Sigma}_{\alpha_0}^+=\Sigma_{\unsure{\alpha_0|\alpha}}^-
\]
% margin note written sideways (labels between braces)
\margin{\unsure{j.e.|i.e.}\diagram{\begin{tikzcd}[ampersand replacement=\&,column sep=0pt,row sep=0.9ex,execute at end picture={
\draw ([xshift=-4.0mm]scTop.west) .. controls +(0.7mm,-1mm) and +(1mm,2mm) .. ([xshift=-4.1mm]scMid.west) .. controls +(-1mm,-.8mm) and +(1mm,.5mm) .. ([xshift=-5mm]scMid.west) .. controls +(1mm,-.5mm) and +(1mm,.8mm) .. ([xshift=-4.1mm,yshift=-1mm]scMid.west) .. controls +(.5mm,-2mm) and +(-.6mm,1mm) .. ([xshift=-4.6mm]scEnd.west);
\draw ([xshift=-.9mm]scMid.west) .. controls +(-.4mm,-1mm) and +(.9mm,.7mm) .. ([xshift=-1.1mm]scE.west) .. controls +(-1mm,-.4mm) and +(.7mm,.1mm) .. ([xshift=-1.8mm]scE.west) .. controls +(.7mm,-.1mm) and +(.8mm,.5mm) .. ([xshift=-1.1mm,yshift=-.5mm]scE.west) .. controls +(.4mm,-1mm) and +(-.5mm,.6mm) .. ([xshift=-.8mm]scEnd.west);
\node[rotate=90,inner sep=0pt] at ([xshift=-2.7mm,yshift=-1mm]scE.west) {$\scriptscriptstyle\Sigma_i^-\,\unsure{\rotatebox{-90}{$!$}|\text{ou}|-}\,\Sigma_i^+$};
}]
|[alias=scTop,anchor=west]| \scriptstyle\Sigma_i^+\ \text{ou}\ \Sigma_i^-\\
|[anchor=west]| \scriptstyle E_i\\
|[alias=scMid,anchor=west]| \scriptstyle\unsure{\Sigma_{i\alpha}|\Sigma_{i,\alpha}}\\
|[alias=scE,anchor=west]| \scriptstyle E_\alpha\\
|[alias=scEnd,anchor=west]| \scriptstyle\Sigma_{\alpha_0}^+
\end{tikzcd}}}
d'autre part $\alpha_0\in K_1$ \unsure{signifi|signifie}
\[
\Sigma_{\alpha_0}^+\not\subset\Sigma_i^+
\]
ce qui \deleted{\unsure{impli|impl.}} \unsure{signifie} aussi que $\Sigma_i^+\cap\Sigma_\alpha^+=\varnothing$, cqfd.

\uline{Corollaire du lemme} Pour une classe

\noindent\begin{tabular}{@{}|p{.96\linewidth}@{}}
$V$ de $\Pi_I^{\unsure{\circ|0}}$, son expression
\[
V=X\setminus\bigcup_{\alpha\in K}\widehat{\Sigma}_\alpha^+
\]
du lemme précédent est \uline{unique} -- on a
{\small
\[
\left\{\begin{aligned}
&K=\{\alpha\in I\mid\overline V\cap E_\alpha\ne\varnothing\quad\text{i.e.}\quad\deleted{\unsure{V}}E_\alpha\cap\dot V\ne\varnothing\}\\
&\Sigma_\alpha^+\ \text{est l'unique comp. de }X\setminus\Sigma_\alpha\ \text{contenue dans }X\setminus V
\end{aligned}\right.
\]}
\end{tabular}
% archive sheet number: 20
\scanpage{22}
\deleted{On procède par récurrence sur $\operatorname{card} K$,} \inserted{On note d'abord que} \deleted{pour prouver que} (si $V=\underset{\substack{\rotatebox{90}{$=$}\\ \inserted{\bigcap_{\alpha\in K}\Sigma_\alpha^-}}}{\underbrace{X\setminus\bigcup_{\alpha\in K}\widehat{\Sigma}_\alpha^+}}$, les $\widehat{\Sigma}_\alpha^+$ \unsure{(}$\alpha\in K$ étant \unsure{disjoints|disjointes})
\[
\left\{\begin{aligned}
&\overline V\cap E_\alpha\ne\varnothing\\
&\makebox[0pt][r]{$\longrightarrow$\hspace{.15em}}V\cap\Sigma_\alpha^-\deleted{\unsure{\ne|=}\varnothing}\ne\varnothing
\qquad\text{alors que } V\cap\Sigma_\alpha^+=\varnothing
\end{aligned}\right.
\]
% The marginal arrow crosses the brace and ends at V in its second row.
\reversemarginpar\margin{\raggedleft\raisebox{.8\baselineskip}{trivial}}\normalmarginpar

De façon précise, je dis que
\[
\begin{aligned}
\overline V\cap E_\alpha & = \overline{\Sigma_\alpha^-}\cap E_\alpha\qquad\text{(qui est donc }\ne\varnothing\text{)}\\[-2pt]
\underset{\inserted{\overline V\cap\widehat{\Sigma}_\alpha^+}}{\rotatebox{90}{$=$}}&
\hspace{1em}\underset{\inserted{\overline{\Sigma_\alpha^-}\cap\widehat{\Sigma}_\alpha^+}}{\rotatebox{90}{$\unsure{\deleted{=}|=|\Rightarrow}$}}
\end{aligned}
\]

\deleted{Comme $V\subset\Sigma_\alpha^-$, on a $V\cap E_\alpha\subset\overline{\Sigma_\alpha^-}\cap E_\alpha$, prouvons l'inclusion inverse}
\[
\deleted{\overline{\Sigma_\alpha^-}\cap E_\alpha\subset\overline V\cap E_\alpha}
\]
\deleted{Soit $x\in\overline{\Sigma_\alpha^-}\cap E_\alpha$, prouvons $x\in\overline V\cap E_\alpha$ i.e. $x\in\overline V$, soit $W$ un voisinage de $x$, montrons qu'il rencontre}

\deleted{Or} \deleted{\unsure{on}} \deleted{peut considérer cette} \deleted{\unsure{formule|relation}} \deleted{\inserted{égalité} comme égalité dans l'espace}
\[
\deleted{\unsure{Y|Z}=\widehat{\Sigma}_\alpha^-=\Sigma_\alpha^-\cup E_\alpha}
\]
\deleted{dans lequel on a deux parties fermées disjointes}

Or dans $X$ on a les deux parties \uline{fermées} \uline{disjointes}
\[
Y=\widehat{\Sigma}_\alpha^+,\qquad Z=\bigcup_{\beta\in K\setminus\{\alpha\}}\widehat{\Sigma}_\beta^+
\]
et on doit prouver
\[
\overline{(X\setminus(Y\cup Z))}\cap Y=\overline{(X\setminus Y)}\cap Y
\]
ce qui est immédiat
% archive sheet number: 21
\scanpage{23}
Soit maintenant $\beta\in I\setminus K$, prouvons que l'on a $\overline V\cap E_\beta=\varnothing$. En effet, $V$ ne \deleted{contient} \inserted{rencontre} aucun $E_i$ $(i\in I)$ donc $E_\beta\subset X\setminus V=\mathop{\unsure{\deleted{\bigvee}\bigcup|\bigcup}}\limits_{\alpha\in K}\widehat{\deleted{\unsure{I|E}}\Sigma}_\alpha^+$, \deleted{\unsure{donc si}} je dis que si $\beta\in I\setminus K$, on a en fait
\[
E_\beta\subset\bigcup\Sigma_\alpha^+\qquad\text{(ouvert disjoint de }V\text{)}
\]
% slanted
\reversemarginpar\margin{\makebox[\linewidth][r]{\rotatebox{24}{\shortstack[l]{C'est trivial, \deleted{\illegible{1}}\\du fait que les $E_i$ $(i\in\unsure{I|J})$\\sont disjoints}}}}\normalmarginpar
ce qui \unsure{prouvera} l'assertion -- mieux, je dis que $\exists\alpha\in K$ tel que $E_\beta\subset\Sigma_\alpha^+$. Cela revient \unsure{au}

\uline{Lemme} Soit $K\subset I$, \deleted{tel que les} \deleted{et} des $\Sigma_\alpha^+$ $(\alpha\in K)$ \unsure{tels} que
\[
\alpha\mathrel{\unsure{\ne|\deleted{\ne}{,}}}\beta\in K,\ \alpha\ne\beta\ \Rightarrow\ \Sigma_\alpha^+\cap\Sigma_\beta^+=\varnothing.
\]
Soit $i\in I\setminus K$ tel que $E_i\subset\bigcup\Sigma_\alpha^+$, alors \deleted{$E_i$ est \unsure{contenu|contenue}} \inserted{$\exists\alpha\in K$ tel que} $E_i\subset\Sigma_\alpha^+$.

Cela \deleted{signifie que} \unsure{provient|précise|précise le} \unsure{du|des|sens}

\uline{Corollaire} Avec les hyp. \unsure{préc.} sur $(\Sigma_i^\alpha)_{\alpha\in\unsure{K|I}}$,

\noindent\begin{tabular}{@{}|p{.96\linewidth}@{}}
la partition $\Pi_K^{\unsure{\circ|0}}$ est formée des $\Sigma_\alpha^+$ $(\alpha\in K)$ et de $V=X\setminus\mathop{\unsure{\rlap{\raisebox{.65ex}{$\scriptstyle-$}}\bigcup|\bigcup}}\limits_{\alpha\in K}\widehat{\Sigma}_\alpha^+$.
\end{tabular}

Immédiat par récurrence sur $\operatorname{card} K\geq 0$.
% archive sheet number: 22
\scanpage{24}
Considérons la partition $\Pi_I^{\unsure{\circ|0}}$ de $X\setminus\bigcup_{\unsure{\alpha'|i}\in I}E_{\alpha'}$, soit $S$ l'\unsure{ens|ens.} des classes, notées $U_s$ $(s\in S)$. \deleted{\unsure{On a}} Considérons la \deleted{\uline{Lemme} $\forall i\in I$,} \inserted{Relation $R$ entre $S$ et $I$}
\[
\deleted{\unsure{U|E}}(s,i)\in R\ \overset{\text{\unsure{def|déf}}}{\Longleftrightarrow}\ \overline{U_s}\cap E_i\ne\varnothing.
\]

\uline{Lemme} Cette relation fait de $(S,I\unsure{)\deleted{R})|\deleted{)}{,}R)}$

\noindent\begin{tabular}{@{}|p{.94\linewidth}@{}}
un graphe strict ayant $S$ comme \unsure{ens des|ens de} sommets, $I$ comme ens. d'\unsure{arêtes|arètes|aretes}.) \deleted{Cette relation} \unsure{Le|Ce} graphe est un \uline{arbre}
\end{tabular}

\unsure{le|Le} fait qu'on a un graphe strict signifie que a) $\forall i\in I$, Il \deleted{\unsure{existe|ex}} y a exactement deux $U_s$ tels que $\overline{U_s}\cap E_i\ne\varnothing$, et b) $i$ est caractérisé par l'ens de ces deux $U_s$.

\unsure{Soient|Soit} $\Sigma_i^+,\Sigma_i^-$ les deux comp. \unsure{connexes} de $X\setminus E_i$, alors a) se précise en : il y a \unsure{exactement|exactemt} \uline{un} $U_{s_{\unsure{+|-}}}$ contenu dans $\Sigma_i^{\unsure{+|-}}$ tel que $\overline{U_{s_+}}\cap E_i\ne\varnothing$ (et de $=$ pour $\Sigma_i^-$, \unsure{dé|qui} définit un $U_{s_-}\subset\Sigma_i^-$ tel que $\overline{U_{s_-}}\cap E_i\ne\varnothing$),
% archive sheet number: 23
\scanpage{25}
\deleted{Soit en effet} Si on a $U_s\subset\deleted{\unsure{E|I}}\Sigma_i^-$, $\overline{U_s}\deleted{\unsure{=}}\cap E_{\unsure{i|j}}\ne\varnothing$, \deleted{et si} alors on a
\[
\begin{aligned}
U_s=X\setminus\mathop{\unsure{\rlap{\raisebox{.65ex}{$\scriptstyle-$}}\bigcup|\deleted{\bigcap}\bigcup}}\limits_{\alpha\in K}\widehat{\Sigma}_\alpha^+
&\qquad\text{où}\quad i\in K\subset I\\
&\forall\alpha\in K\qquad\Sigma_\alpha^+\in\pi_0(X\setminus E_\alpha)\ \deleted{\unsure{\text{les}|K}}\\
&\deleted{\unsure{U\Sigma|\bigcup\Sigma|U}}\ \text{les }\Sigma_\alpha^+\ \text{\unsure{mut\textsuperscript{t}.|mut.} disjoints}
\end{aligned}
\]
Il faut voir que $\deleted{K}\inserted{K\setminus\{i\}}$ et les $\Sigma_\alpha^+$ $(\alpha\in K\setminus\{i\})$ sont \unsure{determinés|déterminés|determines} en termes de $i$ et de $\Sigma_i^+$ i.e. $\Sigma_i^-$. Je dis que
\[
K\setminus\{i\}=\bigl\{\alpha\in I\mid E_\alpha\subset\Sigma_\alpha^-
\]
% archive sheet number: 24
\scanpage{26}
Soit $X$ un ens. Une \uline{décomposition binaire} de $X$ est un ens de deux parties de $X$ $(\Sigma_{\alpha}^{\alpha})_{\alpha\in\unsure{\mathcal E|\mathscr E|\varepsilon}}$ (les ``hémisphères'') dont la réunion est $X$, et dont \deleted{l'intersection} aucune n'est contenue dans l'autre -- on désigne par $E$ (``équateur'') l'intersection (on n'exclut pas qu'elle soit vide). Quand $X$ est un espace topologique, on s'intéresse au cas où les hémisphères $\Sigma_{\alpha}$ sont fermés, (et \inserted{de} plus, le plus souvent, \inserted{$\Sigma_{\alpha}^{\alpha}$ et $\Sigma_{\deleted{\unsure{\alpha}}}^{\alpha}\setminus E=\Sigma_{\deleted{\unsure{\alpha}}}^{\deleted{\unsure{\alpha}}\alpha,\unsure{0|\circ|o}}$} connexes \deleted{\unsure{$\Sigma_{\alpha}$|les}} et où $E\ne\varnothing$); si les $\Sigma_{\deleted{\unsure{\alpha}}}^{\alpha}$ connexes et $E\mathrel{\unsure{\ne|+}}\varnothing$, $X$ est connexe, si \deleted{les $\Sigma_{\alpha}^{\alpha,0}$ connexes}, \inserted{$X$ connexe, \deleted{et}} on conclut \deleted{\unsure{aussitôt}} que $\overline{\Sigma^{\alpha,0}}\cap E\ne\varnothing$
% margin note written sideways
\reversemarginpar\margin{\diagram{\rotatebox{52}{\begin{tikzcd}[ampersand replacement=\&,column sep=2.8em,execute at end picture={
\node[anchor=base west,inner sep=0pt] at ([xshift=1.4mm,yshift=4.8mm]scLeft.center) {\scriptsize diagramme};
}]
|[alias=scLeft]| \bullet\arrow[r,dash,"{\deleted{\unsure{E|\Sigma|B}}}","{\scriptstyle\Sigma^{-}\supset E_{\unsure{\alpha|{}}}\subset\Sigma^{+}}"'] \& \bullet
\end{tikzcd}}}}\normalmarginpar%
(car $\overline{\Sigma^{\alpha,0}}\cap E=\overline{\Sigma^{\alpha,0}}\setminus\Sigma^{\alpha,0}$, et $\Sigma^{\alpha,0}$ ne peut être à la fois ouvert et fermé car $\Sigma^{\unsure{-,0|\alpha,0|-\alpha,0}}$, est un ouvert $\ne\varnothing$ et $\ne X$)

Considérons deux décompositions binaires $\{\Sigma_1^{+},\Sigma_1^{-}\},\unsure{(|[}\Sigma_2^{+},\Sigma_2^{-}\}$ \deleted{\unsure{distinctes}} \inserted{de $X$,} on dit qu'elles sont \uline{non \unsure{panachées|penchées|parallèles}} si \inserted{elles sont \uline{distinctes}, et si} on peut trouver $\Sigma_1^{+}$ dans l'une, $\Sigma_2^{+}$ dans l'autre, tels que
\[
\Sigma_1^{+}\cap\Sigma_2^{+}=\varnothing\qquad(\text{d'où }E_1\cap E_2=\varnothing)
\]
on a alors
\[
\begin{aligned}
\Sigma_1^{+}&\mathrel{\unsure{\deleted{\cap}\subset|\subset}}\Sigma_2^{-}\ \deleted{\unsure{\ne\varnothing}}\quad\text{i.e.}\quad\Sigma_1^{+}\cap\Sigma_2^{-}=\Sigma_1^{+}\\
\Sigma_2^{+}&\subset\Sigma_1^{-}\qquad\qquad\text{\unsure{i.e.|ie}}\quad\Sigma_2^{+}\cap\Sigma_1^{-}=\Sigma_2^{+}\\
\Sigma_{12}&\overset{\mathrm{def}}{=}\Sigma_1^{-}\cap\Sigma_2^{-}\supset E_1\cup E_2,\quad\Sigma_1^{-}\cap\Sigma_2^{-}\ne\varnothing
\end{aligned}
\]
% archive sheet number: 25
